\section*{第 \S\arabic{exercisegroup} 节习题}
\phantomsection
\addcontentsline{toc}{section}{第 \protect\S\arabic{exercisegroup} 节习题}

\begin{enumerate}
\def\labelenumi{\arabic{enumi})}
\item
  设 \(f\) 是定义在开区间 \(I = ]a, b[\)（含于 \(R\)）上且左连续的实函数；假设在 \(B\)（\(I\) 中 \(I\) 的一个可数子集的余集）的所有点上，函数 \(f\) 向右递增，也就是说，在每个点 \(x \in B\) 处存在 \(y\)，使得 \(x < y \leqslant b\) 并且对一切 \(z\)（满足 \(x \leqslant z < y\)）有 \(f(x) \leqslant f(z)\)。证明 \(f\) 在 \(I\) 上递增（按命题 2 的方式论证）。
\item
  在 p 进数域 \(Q_{p}\)（Gen.~Top., III, p.~322, 习题 23）中，每个 p 进整数 \(x \in Z_{p}\) 都具有形如 \(x = a_{0} + a_{1}p + \cdots + a_{n}p^{n} + \cdots\) 的唯一展开，其中诸 \(a_{j}\) 是使对每个 j 有 \(0 \leqslant a_{j} \leqslant p - 1\) 的有理整数。对每个 \(z \in Z_{p}\)，令
\end{enumerate}

\[
f (x) = a _ {0} + a _ {1} p ^ {2} + \dots + a _ {n} p ^ {2 n} + \dots ;
\]

证明在 \(Z_{p}\) 上，f 是连续函数，它在任何点的任何邻域上都不取常值，却在每一点都有等于 0 的导数。

\begin{enumerate}
\def\labelenumi{\arabic{enumi})}
\setcounter{enumi}{2}
\tightlist
\item
  \begin{enumerate}
  \def\labelenumii{\alph{enumii})}
  \tightlist
  \item
    设 K 是三分康托集（Gen.~Top., IV, p.~338），设 \(I_{n,p}\) 是 K 的 \(2^{n}\) 个长度为 \(1/3^{n+1}\)（\(1 \leqslant p \leqslant 2^{n}\)）的邻接区间，而 \(K_{n,p}\) 是 \(2^{n+1}\) 个长度为 \(1/3^{n+1}\) 的闭区间，其并是诸 \(I_{m,p}\)（对 \(m \leqslant n\)）之并的余集。设 \(\alpha\) 是使 \(1 < \alpha < 3/2\) 成立的数；对每个 n，用 \(f_{n}\) 表示 {[}0, 1{]} 上的连续递增函数，它在 x = 0 处等于 0，在每个区间 \(I_{m,p}\)（对 \(m \leqslant n\)）上取常值，在每个区间 \(K_{n,p}\)（\(1 \leqslant p \leqslant 2^{n+1}\)）上是仿射线性的，并且在这些后一区间的每一个的内部满足 \(f_{d}'(x) = \alpha^{n+1}\)。证明通项为 \(f_{n}\) 的级数在 {[}0, 1{]} 上一致收敛，其和是一个函数 f，它在 {[}0, 1{[} 中处处具有右导数（有限或否），并且 \(f_{r}'(x) = +\infty\) 在 K 的每个异于诸邻接区间 \(I_{n,p}\) 之左端点的点处成立。
  \end{enumerate}
\end{enumerate}

\begin{enumerate}
\def\labelenumi{\alph{enumi})}
\setcounter{enumi}{1}
\tightlist
\item
  设 g 是 \([0, 1]\) 到其自身上的连续递增映射，在每个区间 \(I_{n,p}\) 上取常值（Gen.~Top., IV, p.~403, 习题 9）。若 \(h = f + g\)，证明 h 具有等于 \(f_{d}^{\prime}(x)\) 的右导数（在 \([0, 1]\) 的每个点 x 处）。
\end{enumerate}

\begin{enumerate}
\def\labelenumi{\arabic{enumi})}
\setcounter{enumi}{3}
\tightlist
\item
  设 \(f\) 是实值有限函数，在紧区间 \([a, b]\)（含于 \(\mathbf{R}\)）上连续，并且在开区间 \(]a, b[\) 的每一点都有右导数。设 \(m\) 与 \(M\) 是 \(f_d'\) 在 \(]a, b[\) 上的下确界与上确界（有限或否）。
\end{enumerate}

\begin{enumerate}
\def\labelenumi{\alph{enumi})}
\item
  证明当 x 与 y 跑遍 {]}a, b{[} 而保持 \(x \neq y\) 时，\((f(x) - f(y))/(x - y)\) 的值所成的集包含 {]}m, M{[} 且含于 {[}m, M{]}。（归结为证明：若 \(f_{d}^{\prime}\) 在 {]}a, b{[} 的两点 c, d 处取异号的两个值（其中 c \textless{} d），则存在区间 {]}c, d{[} 的两个不同的点，使 f 在其上取同一个值。）
\item
  若进一步 f 在 \(]a\) , b{[} 的每一点都有左导数，则 \(f_{d}^{\prime}\) 与 \(f_{g}^{\prime}\) 在 \(]a\) , b{[} 上的下确界（相应地，上确界）相等。
\item
  由此推出：若 \(f\) 在 \(]a, b[\) 上可导，则 \(f'\) 下每个含于 \(]a, b[\) 的区间的像本身是一个区间，从而是连通的（利用 \(a\)）。
\end{enumerate}

\begin{enumerate}
\def\labelenumi{\arabic{enumi})}
\setcounter{enumi}{4}
\item
  设 \(\mathbf{f}\) 是如下定义的从 \(I = [0, 1]\) 到 \(\mathbf{R}^3\) 的向量映射：对 \(0 \leqslant t \leqslant \frac{1}{4}\)，\(\mathbf{f}(t) = (-4t, 0, 0)\)；对 \(\frac{1}{4} \leqslant t \leqslant \frac{1}{2}\)，令 \(\mathbf{f}(t) = (-1, 4t - 1, 0)\)；对 \(\frac{1}{2} \leqslant t \leqslant \frac{3}{4}\)，令 \(\mathbf{f}(t) = (-1, 1, 4t - 2)\)；最后，对 \(\frac{3}{4} \leqslant t \leqslant 1\)，取 \(\mathbf{f}(t) = (4t - 4, 1, 1)\)。证明由集 \(\mathbf{f}_d'(\mathrm{I})\) 生成的凸集不同于 \(\frac{\mathbf{f}(y) - \mathbf{f}(x)}{y - x}\) 的值集（当 \((x, y)\) 跑遍 \(I\) 的互异点对所成之集时）的闭包（参见~习题 \(4a\)））。
\item
  在区间 \(\mathbf{I} = [-1, + 1]\) 上考虑向量函数 \(\mathbf{f}\)，取值于 \(\mathbf{R}^2\)，定义如下：\(\mathbf{f}(t) = (0,0)\)（当 \(-1\leqslant t\leqslant 0\) 时）；
\end{enumerate}

\[
\mathbf {f} (t) = \left(t ^ {2} \sin \frac {1}{t}, t ^ {2} \cos \frac {1}{t}\right)
\]

当 \(0 \leqslant t \leqslant 1\) 时。证明 \(\mathbf{f}\) 在 \([-1, +1[\) 上可导，但这个区间在 \(\mathbf{f}'\) 下的像不是 \(\mathbf{R}^2\) 中的连通集（\(cf.\) 习题 \(4c\)））。

\begin{enumerate}
\def\labelenumi{\arabic{enumi})}
\setcounter{enumi}{6}
\item
  设 f 是定义在开区间 \(I \subset R\) 上、取值于 R 上的赋范空间 E 中的连续向量函数，并且在 I 的每一点都具有右导数。证明 I 中使 f 具有导数的点所成的集是 I 的一个瘦子集的余集（利用 I, p.~36 的习题 5 b) 以及 I, p.~18 的命题 6）。
\item
  在区间 {[}0, 1{]} 上考虑一个两两不相交的开区间族 \((\mathrm{I}_{n,p})\)，归纳地定义如下：整数 \(n\) 取遍一切值 \(\geqslant 0\)；对 \(n\) 的每个值，整数 \(p\) 取值 \(1,2,\ldots ,2^n\)；有 \(\mathrm{I}_{0,1} = ]\frac{1}{3},\frac{2}{3}[\)；若 \(\mathrm{J}_n\) 是区间 \(\mathrm{I}_{m,p}\)（与诸数 \(m\leqslant n\) 相应）之并，则 \(\mathrm{J}_n\) 的余集是 \(2^{n + 1}\) 个两两不相交的闭区间 \(\mathrm{K}_{n,p}\)（\(1\leqslant p\leqslant 2^{n + 1}\)）之并。若 \(\mathrm{K}_{n,p}\) 是一个区间 \([a,b]\)，则取 \(\mathrm{I}_{n + 1,p}\) 为以 \(b - \frac{b - a}{3}\left(1 + \frac{1}{2^n}\right)\) 与 \(b - \frac{b - a}{3.2^n}\) 为端点的开区间。设 E 是完备集，即 {[}0, 1{]} 相对于诸 \(\mathrm{I}_{n,p}\) 之并的余集。在 {[}0, 1{]} 上定义一个连续实函数 \(f\)，它在 {[}0, 1{]} 的每一点都具有右导数，但在 E 的由异于 E 的邻接区间的端点的点组成的不可数子集的每一点处没有左导数（参见~习题 7）。（在 E 上取 \(f(x) = 0\)，并把 \(f\) 在每个区间 \(\mathrm{I}_{n,p}\) 上适当地加以定义，使得对每个 \(x\in \mathrm{E}\)，都存在点 \(y < x\)（不属于 E，任意接近 \(x\)），并且满足 \(\frac{f(y) - f(x)}{y - x} = -1\)。）
\item
  设 \(f\) 与 \(g\) 是两个实值有限函数，在 \([a, b]\) 上连续，都在 \(]a, b[\) 上有有限导数；证明存在 \(c\)，使得 \(a < c < b\)，并且
\end{enumerate}

\[
\left| \begin{array}{c c} f (b) - f (a) & g (b) - g (a) \\ f ^ {\prime} (c) & g ^ {\prime} (c) \end{array} \right| = 0.
\]

¶ 10) 设 \(f\) 与 \(g\) 是两个实值有限函数，严格正、连续且在开区间 I 上可导。证明若 \(f'\) 与 \(g'\) 严格正且 \(f'/g'\) 在 I 上严格递增，则或者 \(f/g\) 在 I 上严格递增，或者存在数 \(c \in I\)，使得 \(f / g\) 对 \(x \leqslant c\) 严格递减而对 \(x \geqslant c\) 严格递增（注意若有 \(f'(x) / g'(x) < f(x) / g(x)\) 则也有

\[
f ^ {\prime} (y) / g ^ {\prime} (y) <   f (y) / g (y)
\]

对一切 y \textless{} x 成立）。

\begin{enumerate}
\def\labelenumi{\arabic{enumi})}
\setcounter{enumi}{10}
\tightlist
\item
  设 \(f\) 是复函数，在开区间 \(I\) 上连续、处处不为零，并且在 \(I\) 的每一点都具有右导数。为使 \(|f|\) 在 \(I\) 上递增，必须且只需 \(\mathcal{R}(f_d' / f) \geqslant 0\) 在 \(I\) 上成立。
\end{enumerate}

¶ 12) 设 \(f\) 是开区间 I 上的可导实函数，\(g\) 是它在 I 上的导数，\([a, b]\) 是含于 I 的紧区间；假设 \(g\) 在开区间 \([a, b]\) 上可导，但在点 \(a\)（相应地，\(b\)）处未必右（相应地，左）连续；证明存在 \(c\)，使得 \(a < c < b\)，并且

\[
g (b) - g (a) = (b - a) g ^ {\prime} (c)
\]

（利用 I, p.~36 的习题 4 c)）。

\begin{enumerate}
\def\labelenumi{\arabic{enumi})}
\setcounter{enumi}{12}
\tightlist
\item
  我们把向量函数 \(\mathbf{f}\) 在一点 \(x_0\)（内于 \(\mathbf{f}\) 的定义区间）处的对称导数定义为：\(\frac{\mathbf{f}(x_0 + h) - \mathbf{f}(x_0 - h)}{2h}\) 当 \(h\) 保持 \(>0\) 而趋于 0 时的极限（当它存在时）。
\end{enumerate}

\begin{enumerate}
\def\labelenumi{\alph{enumi})}
\item
  把 § 1 中对导数建立的运算法则推广到对称导数。
\item
  证明当把“右导数”一词换成“对称导数”时，§ 2 的定理 1 与定理 2 仍然成立。
\end{enumerate}

\begin{enumerate}
\def\labelenumi{\arabic{enumi})}
\setcounter{enumi}{13}
\tightlist
\item
  设 \(\mathbf{f}\) 是定义在紧区间 \(\mathrm{I} = [a, b]\)（含于 \(\mathbf{R}\)）上且连续的向量函数，取值于 \(\mathbf{R}\) 上的一个赋范空间中。假设 \(\mathbf{f}\) 在此区间的某个可数子集 A 相对于 \([a, b[\) 的余集的所有点处都具有右导数。证明存在一点 \(x \in ]a\)、\(b[\cap ∁A\)，使得
\end{enumerate}

\[
\left\| \mathbf {f} (b) - \mathbf {f} (a) \right\| \leqslant \left\| \mathbf {f} _ {d} ^ {\prime} (x) \right\| (b - a).
\]

（用反证法论证，把 \([a, b[\) 分解为三个区间 \([a, t[, [t, t + h[\) 与 \([t + h, b[\)，满足 \(t \notin A\)；若 \(k = \| \mathbf{f}(b) - \mathbf{f}(a)\| / (b - a)\)，则注意当 \(h\) 充分小时有 \(\| \mathbf{f}(t + h) - \mathbf{f}(t)\| < k.h\)，并对其余区间使用 I, p.~15 的定理 2。）

\setcounter{exercisegroup}{2}
\refstepcounter{exercisegroup}
